\documentclass[11pt,a4paper]{article}
\usepackage[utf8]{inputenc}
\usepackage[T1]{fontenc}
\usepackage[margin=2.2cm]{geometry}
\usepackage{booktabs,hyperref,listings,xcolor,array}
\definecolor{academyblue}{HTML}{00459A}
\hypersetup{colorlinks=true,linkcolor=academyblue,urlcolor=academyblue}
\title{M301 Software Engineering Final Project\\Async Research Assistant}
\author{Samir Həsənov (@Hasanov57) \\
Elmir Əsgərov (@elmiresgerov) \\
Ramil Məmmədəliyev (@ramilmammadaliyev)}
\date{September 18, 2026}

\begin{document}
\maketitle

\begin{center}
\textbf{Repository:} \url{https://github.com/Hasanov57/AI-Academy---SE---Final-Project---AI-Researcher}\\
\textbf{Final tag:} \texttt{v1.0-final}
\end{center}

\section*{Executive Summary}
We built a command-line research assistant that queries Wikipedia, arXiv and Tavily concurrently
and produces a concise answer with numeric citations. The student-owned layer wraps the supplied
\texttt{ai/} package without modifying its interface. It adds typed configuration, PostgreSQL
persistence, a TTL source cache, retries, timeouts, rate limiting and structured logging. A bounded
\texttt{asyncio} pipeline isolates failed sources so remaining evidence can still reach synthesis.
The live benchmark measured a 1.440-second sequential median and a 0.544-second concurrent
median, a 2.64x speedup with application caching disabled, while sixty-four offline tests pass with
73.56 percent coverage. A real Gemini request incompatibility taught us to keep provider
adaptations outside the supplied package; if starting again, we would enforce that boundary from
the first commit and add semantic citation-entailment checks.

\tableofcontents

\section{Project Overview}
\subsection{Problem framing and scope}
A user submits one question. The application gathers short, relevant excerpts from three distinct
origins and passes only that evidence to the provided synthesizer. The output contains an answer,
inline \texttt{[N]} markers, a reference list, source timings and any degradation warnings. The
mandatory interface is \texttt{python -m researcher ask "question"}; \texttt{--sources} restricts
origins and \texttt{--no-cache} bypasses cached reads and writes. The optional Streamlit interface
reuses the same core use case for project-day interaction while the CLI remains the required entry point.

\subsection{Provider choice and AI contract}
Google Gemini \texttt{gemini-3.5-flash-lite} performs synthesis and Tavily supplies general web results. Wikipedia
and arXiv need no credentials. All calls cross the provided \texttt{ai.fetch\_*} and
\texttt{ai.synthesize} boundary. Our code instantiates the supplied provider adapters but never
calls an upstream SDK directly.

\section{Architecture}
\subsection{Module map}
The CLI creates a \texttt{ResearchRequest}; \texttt{Researcher} coordinates source retrieval,
synthesis and persistence. \texttt{SourceOrchestrator} owns bounded task creation.
\texttt{SourceService} owns cache, retry, timeout, rate-limit and logging policy.
\texttt{SynthesisService} moves the blocking adapter to a worker thread and validates citations.
Both \texttt{PostgresRepository} and \texttt{MemoryRepository} implement our
\texttt{ResearchRepository} abstract base class.

\begin{verbatim}
CLI -> Researcher -> SourceOrchestrator -> SourceService -> provided ai.fetch_*
                 |                         |
                 |                         +-> PostgreSQL TTL cache
                 +-> SynthesisService -> provided ai.synthesize
                 +-> PostgreSQL session, source and citation audit
\end{verbatim}

\subsection{OOP design and data flow}
Stateful resources use classes: repository pools, rate limiters, wrappers and the orchestrator.
Pure transformations such as query canonicalization and URL de-duplication remain functions.
Pydantic models cross module boundaries; raw dictionaries remain limited to serialization edges.
Constructor injection lets tests replace source callables and the LLM without network access.

\section{Concurrency and Performance}
The three source calls are independent and I/O-bound, so \texttt{asyncio} is more suitable than
multiprocessing. One shared \texttt{httpx.AsyncClient} reuses connections. A configurable semaphore
bounds parallelism, while \texttt{asyncio.timeout} applies a separate deadline to each origin.
Synthesis uses \texttt{asyncio.to\_thread} because the supplied LLM interface is synchronous.

\begin{center}
\begin{tabular}{lrrr}
\toprule
Workload & Sequential & Concurrent & Speedup\\
\midrule
Three live sources, cache bypassed & 1.440 s & 0.544 s & 2.64x\\
\bottomrule
\end{tabular}
\end{center}

The measurement used three live runs on the same Windows machine with application cache disabled.
Sequential times were 1.440, 1.439 and 1.523 seconds; concurrent times were 0.535, 0.544 and
0.677 seconds. Retrieval, not LLM synthesis, was measured so source parallelism was isolated. The
new bottleneck is the slowest provider plus network variance. arXiv receives at most one new
request per second from a process.

\section{Robustness and Error Handling}
Transient provider, HTTP and timeout errors receive up to three attempts with exponential backoff.
Every call has a deadline and emits start, finish, timing, count and failure logs without secrets.
The orchestrator converts a failed source into a typed outcome and warning. If every source fails,
the system stops before synthesis and prints a clean error.

During live testing Gemini returned HTTP 400 when a provider-specific thinking option was not
accepted by the selected model. The CLI displayed a clean provider error and the logs identified
the model and status without exposing the key. We reverted the supplied adapter to its original
course version and introduced a student-owned \texttt{ResearchLLM} wrapper, preserving the
instructor contract. Separate tests force source failure, repeated timeouts and invalid citations.

Questions are stripped, whitespace-normalized, limited to 1,000 characters and checked for at least
one supported source. Citation validation requires at least one marker, rejects out-of-range markers,
requires marker/reference equality and checks each substantive sentence for a citation.

\section{Storage and Caching}
PostgreSQL stores \texttt{research\_sessions}, ordered \texttt{research\_sources}, per-origin
\texttt{source\_attempts}, \texttt{citations} and \texttt{source\_cache}. Foreign keys cascade from
the session, and indexes cover cache expiry and session creation time. One transaction saves a
complete result.

The cache key is the source selector plus a case-folded, whitespace-normalized query. Source lists
are stored as JSONB with an expiration timestamp. The default TTL is 24 hours. Expired entries are
removed lazily on access, and \texttt{--no-cache} bypasses both cache reads and writes.

\section{Testing}
The suite contains 64 offline tests and reaches 73.56 percent branch-aware coverage of the
student-owned package. The instructor smoke tests remain unchanged. Tests cover cache hit and
expiry, transient retry, timeout exhaustion, citation rejection, input validation, URL
de-duplication, graceful degradation and the end-to-end use case. The concurrency test records three
simultaneously active source calls and compares its wall clock with the sequential baseline.
\texttt{ruff} and strict \texttt{mypy} both complete without issues.

\section{Deployment and Reproducibility}
The last verified single-stage image was 515.4 MB. It uses Python 3.12 slim, installs pinned dependencies, copies the
provided and student packages, drops root privileges and runs the required command. Docker Compose
supplies PostgreSQL 17.6 and waits for its health check. Secrets remain in an ignored
\texttt{.env}; \texttt{.env.example} contains only empty fields and non-secret configuration.

\section{Limitations and Future Work}
\begin{itemize}
  \item Citation syntax checks cannot establish semantic entailment.
  \item A single PostgreSQL instance has no high-availability replication.
  \item In-process arXiv limiting does not coordinate across multiple containers.
  \item Upstream excerpts may omit context even when every service is healthy.
  \item Provider failover, streaming answers and public hosting remain outside the first release;
        the included Streamlit UI runs through Docker Compose.
\end{itemize}

\section{Tools and Acknowledgements}
AI tools were used for architecture planning and technical assistance during development. The team
reviewed and adapted the suggestions and verified the final system through the complete offline
suite and live demonstrations. The provided \texttt{ai/} package and course templates came from
AI Academy.

\section*{References}
\begin{enumerate}
  \item AI Academy, \emph{M301 Software Engineering Final Project Brief}, Spring 2026.
  \item Gemini API documentation: \url{https://ai.google.dev/gemini-api/docs}
  \item Tavily API documentation: \url{https://docs.tavily.com/}
  \item MediaWiki API documentation: \url{https://www.mediawiki.org/wiki/API:Main_page}
  \item arXiv API documentation: \url{https://info.arxiv.org/help/api/index.html}
  \item Python asyncio documentation: \url{https://docs.python.org/3/library/asyncio.html}
\end{enumerate}

\appendix
\section{Reproducing the benchmark}
\begin{verbatim}
python -m pip install -r requirements.txt
python -m pip install -e .
docker compose build researcher
docker compose up -d db
docker compose run --rm researcher python scripts/bench.py --runs 3
\end{verbatim}

\end{document}
